\section{Conclusion and future work}\label{sec:conclusion}

Transparency law presupposes that administrations hold an account of how they operate. Often
they do not. We derived one from an administration's own event log: a catalogue of emergent
process types, each with a typical path, reproducible from stored artefacts and bounded by
measured limits rather than asserted ones.

The evaluation is what the catalogue rests on. Permuting cluster membership while holding the
partition's shape fixed destroys 23 points of downstream accuracy, so the grouping is real. All
nineteen feature-derived clusterings refit exactly, so the artefact is auditable. Against the declared
taxonomy, agreement rises sharply once a single catch-all class covering 30.5\% of the concluded cases is set
aside, which locates the official classification's weakness rather than merely asserting it.

Three results constrain what the method may claim, and they were the reason to run the controls.
What per-cluster modelling contributes depends on how it is measured, and each correction
reduces it. Over the full evaluated population it gains 0.013 on action prediction and 0.063 on
routing when the best of five families is selected per cluster, and 0.004 and 0.035 when the
family is held fixed --- against the 0.083 and 0.156 the conventional within-cluster baseline
suggests. Scored instead on the single population both trace alphabets can be evaluated on, it
gains 0.036 on routing and \emph{loses} 0.009 on actions, in every model family. The
choice of clustering algorithm separates the result very little: sixteen clusterings across four
algorithm families fall within 0.0079 once scored on a common set of cases, on an interval that
is nonetheless too wide to call them equivalent. And the internal validity indices that would
ordinarily choose between them do not measure what they are used for: silhouette tracks how much
a partition discards and how coarse it is --- $\rho = 0.84$ with coverage among the clusterings
that assign every case --- and carries no association with accuracy once cluster size is held
fixed. It is a defensible instrument for choosing $k$ and not for predicting downstream utility.

Four directions follow. Widening the corpus beyond one year is the most valuable and is limited
by cost rather than method: at the observed cluster-size distribution, three years would move
the share of cases sitting in clusters large enough for their path to reach the quality band
where fabrication is low --- 52.9\% in the selected catalogue today, against a threshold of
about 150 traces. We do not project that share forward: a longer window does not simply
multiply cluster sizes, because it also relaxes the right-censoring that shapes which
procedures appear at all, and we have no basis for assuming the type mixture is stationary.
Second, and with the clearest operational payoff, richer early-stage features would support type
inference at intake --- the point at which knowing a case's type could change how it is handled
--- which the temporal-dependency result shows is a harder problem than retrospective
assignment. Third, closing the under-generation gap, where paths currently match under half of a
real case's visits. Fourth, a properly powered test of whether internal validity indices predict
downstream utility, which sixteen partitions cannot settle.
